\documentclass[11pt]{article}
\usepackage[margin=1in]{geometry}
\usepackage{booktabs}
\usepackage{natbib}
\usepackage{amsmath}
\usepackage{amssymb}
\usepackage{longtable}
\usepackage{array}
\usepackage{hyperref}
\usepackage{xcolor}
\usepackage{caption}
\usepackage{pdflscape}

\setlength{\parskip}{0.5em}
\setlength{\parindent}{0pt}

\title{Behavioral responses of flood insurance coverage to premium changes:\\
A literature summary}
\author{Research notes}
\date{\today}

\begin{document}
\maketitle

\begin{landscape}
\section*{Summary table}

The table collects headline estimates of how flood insurance \emph{coverage}
(take-up / policies in force) responds to \emph{premiums}. ``Implied elasticity'' is the
ratio of the percentage change in coverage to the percentage change in premium when a
paper does not report an elasticity directly; flagged with $^\dagger$ and discussed in
the corresponding entry. Entries marked with $^\ddagger$ are based on third-party
summaries rather than my own read of the paper and should be verified before citation.

\vspace{1em}

\renewcommand{\arraystretch}{1.3}
\footnotesize
\begin{tabular}{p{3cm} p{4.5cm} p{3.5cm} p{3.5cm} p{6cm}}
\toprule
\textbf{Study} & \textbf{Identification} & \textbf{$\Delta\pi$} & \textbf{$\Delta I$} & \textbf{Elasticity / response} \\
\midrule
\citet{hennighausen2023}
& DiD: pre- vs.\ post-FIRM in SFHA, BW-2012 / HFIAA-2014; NFIP tract panel 2009--18
& $+14.3\%$ avg.\ ($+24\%$ by 2018)
& $-8.2\%$ avg.\ ($-16.9\%$ by 2018)
& Implied elasticity $\approx -0.64$ (log-ratio); $\approx -0.70$ by 2018$^\dagger$ \\
\citet{collier2026}
& DiD: SRL vs.\ non-SRL pre-FIRM Zone-A primary residences, BW-2012; NFIP policy panel 2009--18 (2009 cohort)
& $+25\%$ in Yr 1; $+25\%$/yr thereafter (prices roughly double by 2019)
& $-26\%$ cumulative non-renewal over Yrs 1--2; no further effect after FY2014
& Author-reported first-year elasticity $\approx -1.0$ (p.~20). Multi-year effective elasticity is much smaller in absolute value: non-renewal saturates after FY2014 despite further price increases$^\dagger$ \\
\citet{mulder2024}
& OLS on renewal indicator exploiting Newly Mapped Subsidy phase-out (calendar-year vs.\ within-year timing); NFIP newly-mapped policy panel, 2015--19
& $+15\%$/yr scheduled at second renewal under Newly Mapped Subsidy
& $-0.17$ pp on renewal probability per $+1\%$ premium
& Author-reported price elasticity $\approx 0.18$ (Table~1, col.~4). Mandate-adjusted voluntary elasticity $\approx 0.37$ (App.~F). Strictly a semi-elasticity; treated as an elasticity in the welfare model$^\dagger$ \\
\citet{bradt2021}$^\ddagger$
& Quasi-experimental BW/HFIAA design covering both SFHA (mandate) and non-SFHA (voluntary); NFIP policy data
& BW/HFIAA-induced premium changes
& Take-up declines reported for both segments
& Reported elasticity $\approx -0.25$ to $-0.29$. Useful for separating mandate-driven from voluntary demand \\
\citet{wagner2022}$^\ddagger$
& House-level NFIP premium variation from BW/HFIAA in 20 Atlantic and Gulf Coast states, 2009--17; structural demand model with adaptation
& BW/HFIAA premium increases
& Take-up response in structural model
& Reported demand elasticity $\approx -0.25$ to $-0.30$. Recovers WTP $<2/3$ of expected loss inside SFHA \\
\citet{gourevitch2025}$^\ddagger$
& DiD-style comparison across ZIPs by quartile of RR2.0 premium increase, 2021--23
& RR2.0-induced premium increases (vary by quartile)
& $-11$ to $-39\%$ new policies; $-5$ to $-13\%$ renewals
& Implied elasticities $-1.38$ to $-0.26$ across subgroups; effects largest in lower-income ZIPs$^\dagger$ \\
\citet{ortega2025}$^\ddagger$
& IV using individual NFIP insurance histories, RR2.0 reform
& RR2.0 premium changes
& Increased exit, reduced entry both inside and just outside the SFHA; intensive-margin adjustments in coverage and deductibles
& No clean elasticity reported in available summaries; complements \citet{gourevitch2025} via individual-level data and intensive margin \\
\citet{browne2000}$^\ddagger$
& State-year cross-sectional NFIP aggregates
& Cross-sectional premium variation
& Policies in force vs.\ total insured amount
& Reported elasticities: $-0.32$ (extensive margin, policies) and $-1.22$ (intensive margin, total insured amount). Identification weak --- premiums and quantities both depend on risk and loss history \\
\citet{atreya2015}$^\ddagger$
& 32-year Georgia county panel
& Within-county premium variation
& County-level take-up
& Reported elasticities $-0.14$ to $-0.31$ depending on specification. Pre-BW benchmark with weak identification by modern standards \\
\bottomrule
\end{tabular}
\end{landscape}

\clearpage

\section{\citet{hennighausen2023} -- \textit{Journal of Housing Economics}}

\paragraph{Focus relative to our question.}
The paper's headline contributions are housing-market effects (prices, volumes) of the
2012 Biggert--Waters Act (BW) and 2014 Homeowner Flood Insurance Affordability Act
(HFIAA). Estimates relevant to the premium--coverage margin appear as a ``first stage''
analysis (Section 2.3, Table 1, p.\ 4). The paper does not report an elasticity directly;
the response measure has to be inferred from two reduced-form DiD coefficients.

\paragraph{Setting and reforms.}
BW (2012) and HFIAA (2014) phased out subsidies for properties built before their
community's first Flood Insurance Rate Map (``pre-FIRM'') located inside the Special
Flood Hazard Area (SFHA, the 100-year floodplain). These rules raised pre-FIRM premiums
by roughly 18--25\% per year toward the full-risk rate, while post-FIRM properties were
unaffected.

\paragraph{Identification strategy (insurance margin).}
Two-way fixed-effects DiD on a balanced panel of NFIP policies aggregated to the
\emph{census tract $\times$ year $\times$ pre/post-FIRM} cell, 2009--2018 (445{,}110
observations across $\sim$29{,}000 tracts):
\begin{equation*}
\log(\text{Insurance})_{itk} = \alpha_1 \, \text{PreFIRM}_{ik}\times \text{Post}_t
+ \delta_{ik} + \eta_{it} + \varepsilon_{itk}.
\end{equation*}
Tract-by-pre/post-FIRM ($\delta_{ik}$) and tract-by-year ($\eta_{it}$) fixed effects
absorb cell-level levels and time-varying tract-wide shocks. The dependent variable is
either the average premium or the policy count, in logs. $\alpha_1$ identifies the
differential change for treated (pre-FIRM) policies after 2012. Standard errors are
clustered by county.

\paragraph{Headline estimates (Table 1, p.\ 4).}
\begin{itemize}
\item \textbf{Premium response (col.\ 1):} $\hat{\alpha}_1^{\text{prem}} = 0.134$
(s.e.\ 0.005); equivalently a $\sim$14.3\% relative increase in average pre-FIRM
premiums over the study period.
\item \textbf{Coverage response (col.\ 2):} $\hat{\alpha}_1^{\text{policy}} = -0.086$
(s.e.\ 0.008); $\sim$8.2\% decrease in pre-FIRM policy counts, $\approx$170{,}600
fewer policies relative to the 2012 pre-FIRM stock (footnote 11, p.\ 4).
\item \textbf{Event-study path (Fig.\ A.1, p.\ 15):} the premium effect grows to
$\sim$24\% and the policy effect to $\sim$16.9\% by 2018; pre-trend Wald tests are
flat for policies ($p=0.112$) but \emph{not} for premiums ($p=0.00$).
\end{itemize}

\paragraph{Implied price-responsiveness of coverage (conversion).}
Both reduced-form coefficients are in logs and come from the same DiD treatment, so a
Wald-style ratio gives an aggregate elasticity:
\[
\eta \;\approx\; \frac{\hat{\alpha}_1^{\text{policy}}}{\hat{\alpha}_1^{\text{prem}}}
\;=\; \frac{-0.086}{0.134} \;\approx\; -0.64
\quad\text{(average 2012--18)},\qquad
\eta_{2018} \;\approx\; \frac{-0.185}{0.265} \;\approx\; -0.70.
\]
(For comparison, the percent-change ratio $-0.082/0.143 \approx -0.57$; the difference
is the standard log-vs.-percent gap at these magnitudes.)\\
\textbf{Caveats (flagged $^\dagger$):}
\begin{itemize}
\item This is an aggregate Wald ratio, not an individual-level price elasticity. It
should be read as ``how much does coverage fall when the average treated premium rises
by 1\%, under BW/HFIAA''.
\item The authors note (p.\ 9) that the premium estimate may understate the increase
faced by the marginal dropper if those facing the largest hikes selected out of the
program; this biases $|\eta|$ \emph{downward in magnitude}.
\item Estimate is local to previously subsidized SFHA properties (pre-FIRM, inside
100-year floodplain). External validity to post-FIRM, non-SFHA, or fully unsubsidized
properties is not established.
\end{itemize}

\paragraph{Margins captured.}
\textbf{Extensive margin only} (number of policies in force). The data do not separate
new non-purchase from cancellations, nor do they capture switching to private flood
insurance. The authors argue the latter is a minor concern: private flood insurance was
$\approx$4\% of the market in 2018 and tends to attract \emph{lower-risk} properties,
not subsidized pre-FIRM ones (footnote, p.\ 4, citing Kousky et al.\ 2018).

\paragraph{Threats to identification the authors discuss.}
(i) Anticipation: argued unlikely because pre-FIRM status was hard to observe
ex-ante and the political backlash suggests the rate hike was a surprise.
(ii) Concurrent shocks: Hurricane Sandy (Oct.\ 2012) is the main worry; results are
robust to dropping NY and NJ (Table A.5).
(iii) Map updates changing SFHA status: robust to dropping tracts with $>$5\% SFHA
re-zoning during 2000--19 (Table A.6).
(iv) Pre-trend in premiums: the event study shows the premium pre-trend is
statistically rejected, which softens but does not undo the premium effect; the
\emph{coverage} pre-trend is flat.

\paragraph{Tangential housing-market results (for context, not focal).}
Repeat-sales DiD on ZTRAX 2008--18 (Tables 3--4): pre-FIRM property \emph{prices}
fall 4.2\% and \emph{transaction volumes} fall 2.3\%. The price decline is roughly
the NPV of the premium increase (\$6{,}600 implied vs.\ \$6{,}207 NPV at a 3\% discount
rate over 30 years; Appendix C), consistent with near-full capitalization.

\paragraph{Comparison the authors draw.}
The authors note their coverage results are consistent with negative price elasticities
of NFIP demand reported by \citet{bradt2021} and \citet{wagner2022}, but do not provide
a numeric comparison.

\section{\citet{collier2026} -- Working paper, April 2026}

\paragraph{Focus relative to our question.}
The headline outcome of interest is renewal of flood insurance among an
extremely high-risk, heavily-subsidized subset of NFIP policyholders. The
premium--coverage margin is the central object of the paper, not a
by-product, so the relevant estimates are reported directly. The authors
themselves compute a back-of-the-envelope first-year price elasticity of
demand of approximately $-1$ (p.~20), which we adopt as the headline number
and contextualize below.

\paragraph{Setting and reform.}
Severe Repetitive Loss (SRL) properties are NFIP-designated homes with
either four claims of $\$5{,}000{+}$ or two claims summing to more than the
property's value within ten years. They are 1\% of policies but $\sim$30\%
of NFIP claims. Pre-reform, the median SRL policyholder paid $\$994$ in
premium against an estimated $\$5{,}200$ in expected loss (median load
$0.19$; Table~3, Figure~2). Biggert--Waters (2012) phased out these
subsidies by raising SRL premiums 25\% per year until the actuarial rate
was reached. Crucially, the renewal letter accompanying the first
post-reform bill (Aug.\ 2013, Appendix Fig.~A1) explicitly told SRL
policyholders that they had been subsidized, that premiums would rise 25\%
annually, and that premiums would never exceed the full-risk rate.

\paragraph{Identification strategy.}
DiD on a policy--year panel of the 2009 cohort followed through 2018
($\sim$2.05M policy-year observations on $\sim$476K unique policies, of
which 2{,}050 are SRL). Treated and control are restricted to single-family
primary residences, pre-FIRM rated, in Flood Zone~A, who paid
\emph{identical} premium rates pre-reform; the rate schedules diverge
only after Q4~2013 (Figure~3, Panel~a; Appendix Tables~A2--A3 reproduce the
NFIP rating tables). The preferred specification (Equation~1, Table~4
column~4) is
\[
y_{it} = \!\!\sum_{j=2010,\,j\ne 2012}^{2018}\!\!\gamma_j\,\mathbb{I}(t{=}j)\times \mathbb{I}(\mathrm{SRL}_i)
+ \beta\,\mathbb{I}(\mathrm{SRL}_i) + X_i'\delta + FE_{\text{tract}\times\text{year}} + \varepsilon_{it},
\]
with $y_{it}$ a renewal indicator, FY2012 the reference year, and standard
errors clustered at the policy level. A quarterly version (Equation~2,
Figure~4) is used for timing.

\paragraph{Headline estimates (Table~4, col.~4; Figure~4).}
\begin{itemize}
\item \textbf{Annual DiD coefficients on $\mathbb{I}(\mathrm{SRL})\times\mathbb{I}(t)$:}
$\hat\gamma_{2013} = -0.18$ (s.e.\ 0.02), $\hat\gamma_{2014} = -0.21$
(s.e.\ 0.02). All later years are statistically indistinguishable from zero.
\item \textbf{Cumulative non-renewal due to the reform: 26\%}
(footnote~17, p.~18; computed as $0.77\times 0.74 - (0.77{-}0.18)(0.74{-}0.21)
\approx 0.257$, using the non-SRL renewal rates as counterfactuals).
\item \textbf{Quarterly peak (Figure~4):} SRL policies due to renew in Q1~2014
were $\sim$50~pp less likely to renew. Effects begin in Q4~2012 (after BW
passage but \emph{before} any premium change) and end by Q3~2014.
\item \textbf{Intensive margin: zero detectable response.} Coverage limits
of those who continue to insure do not move with the reform in 2013--14;
small effects appear only in 2017--18 and may be selection-driven (Internet
Appendix~B.6, Table~B5; Figure~B4).
\end{itemize}

\paragraph{Author-reported elasticity and conversion notes.}
The authors write (p.~20): ``a 25\% price increase reduces renewal by 26\%
during the first year of the reform,'' yielding an implied elasticity
$\approx -1.0$. This is roughly $1.5\times$ the magnitude of the
\citet{hennighausen2023} aggregate Wald ratio. Two important qualifications:
\begin{itemize}
\item \textbf{The first-year elasticity does not extrapolate.} Prices keep
rising 25\% per year through 2019 (Figure~3a; SRL premium roughly doubles
relative to control by 2019), but no additional non-renewal occurs after
FY2014. Computed naively over the full 2013--2019 window, the implied
elasticity would be much smaller in absolute value. The authors interpret
this as a kinked, L-shaped demand curve (Internet Appendix~B.2,
Figure~B2): one group drops in anticipation, a second group drops at the
first price increase, and the remaining group is highly inelastic.
\item \textbf{Cumulative-vs.-flow elasticity.}$^\dagger$ The 26\% figure is
a cumulative two-year non-renewal effect against a 25\% \emph{first-year}
price increase. A more conservative read uses the cumulative two-year
price increase ($\sim$56\% above pre-reform on a compounded base): then
$-0.26/0.56 \approx -0.46$. We flag the headline $-1.0$ as the
authors' figure but caution against using it as a generic price elasticity
of NFIP demand outside this specific high-subsidy, high-risk population
and the first year of treatment.
\item \textbf{Survival vs.\ flow.} The Internet Appendix~B.3 linear survival
specification shows the effect is \emph{persistent}: by 2018, 9~pp fewer SRL
than non-SRL properties from the 2009 cohort remain insured (Table~B1,
col.~4), so this is not pure harvesting.
\end{itemize}

\paragraph{Margins and decomposition of non-renewals.}
\textbf{Extensive margin only}, with the further breakdown (Section~4):
$\sim$25\% of reform-induced non-renewals are policy \emph{cancellations}
(typically associated with home sales / relocation; Figure~5b). Home sales
in tracts with SRL properties rise by $\sim$1 per quarter immediately
post-reform and by $\sim$2.5 per quarter by 2018 (Figure~6a); tract-level
prices fall $\sim$10\% (Figure~6b). At end-of-period (2018), only 6.6\% of
existing uninsured SRL homes had been mitigated (Table~5), implying that
the bulk of non-renewers neither relocated nor mitigated --- they simply
went unprotected.

\paragraph{Heterogeneity / mechanisms (Section~5, Internet Appendix~C--D).}
Treatment-intensity DiD (Eq.~D1, Figure~C2) shows: \emph{no} gradient in
non-renewal with pre-reform contract load (rules out adverse selection on
observables); \emph{no} gradient with announced premium-increase magnitude
or with tract-level income (no support for liquidity constraints);
significant gradient with tract-level low-education share. The timing of
non-renewals tracks media coverage of BW (Figure~C3) more tightly than the
price path itself, which the authors offer as suggestive of belief- /
attention-driven responses rather than purely price-driven ones.

\paragraph{Threats to identification the authors discuss.}
(i) Parallel trends: pre-reform $\hat\gamma_{2010}$ and $\hat\gamma_{2011}$
are 0.01 or smaller and statistically zero (Table~4); the quarterly event
study shows nine flat pre-reform quarters (Figure~4).
(ii) Load imbalance between groups: results are essentially unchanged when
the sample is restricted to loads $<$1, $<$0.75, $<$0.50 (Internet
Appendix~B.4, Table~B2): $\hat\gamma_{2013}$, $\hat\gamma_{2014}$ stay in
$[-0.22, -0.18]$.
(iii) Cohort-specific shocks: re-estimating on the 2010 cohort gives
essentially the same coefficients (Internet Appendix~B.5, Tables~B3--B4).
(iv) HFIAA (2014) confound: the HFIAA capped general NFIP rate increases at
18\% but did not roll back SRL increases, and control-group premiums in
fact grew 4--5\%/yr both before and after (Figure~3a), consistent with
broader homeowners-insurance inflation.

\paragraph{Welfare estimate (for context, not focal).}
A sufficient-statistics calibration (Section~6, Table~6) returns a welfare
loss of $\$4{,}200$ per SRL homeowner per year ($\sigma=3$, $\lambda=1.3$,
$y_{\min}=\$4{,}000$); the indifference threshold would require 18\% of SRL
owners to fully mitigate, far above the $\le$4.2\% upper bound from the
property-level data.

\paragraph{Population to which the estimate applies.}
Highly selected: pre-FIRM, single-family, primary-residence, Flood-Zone-A,
SRL-designated properties whose pre-reform premiums were on average
$\sim$1/5 of expected loss. The first-year elasticity of $\approx -1$ is a
local estimate for households told they were heavily subsidized and that
their subsidy was being removed; it is \emph{not} a generic NFIP price
elasticity.

\section{\citet{mulder2024} -- Working paper, September 2024}

\paragraph{Focus relative to our question.}
The paper's headline contribution is a sufficient-statistics welfare model
of flood-risk \emph{misclassification}, combining a price effect, an
information effect, and an adaptation effect. Our question maps onto the
``price effect'' alone (Section~5.1, Table~1, p.~46; Appendix~F).
The estimate is reported directly by the author as a price elasticity, with
a separate mandate-adjusted version that he uses in the welfare model.

\paragraph{Setting and reform.}
The Newly Mapped Subsidy (started 2015) lets homes \emph{remapped} into the
SFHA pay non-floodplain premiums for two calendar years, after which
premiums increase $\sim$15\% per year toward the floodplain rate. The
phase-out is keyed to the \emph{calendar year} of the new map, not the
month, which is the source of identifying variation.

\paragraph{Identification strategy.}
Linear probability model on a panel of newly-mapped policies, 2015--2019:
\[
\text{Renew}_{iqt} = \alpha_{qt} + \beta X_{iqt} + \sigma\,\Delta P_{iqt} + \varepsilon_{iqt},
\]
where Renew is a 0/1 indicator, $\Delta P_{iqt}$ is the scheduled \%
premium change at renewal under the Newly Mapped subsidy schedule,
$\alpha_{qt}$ are first-policy-year~$\times$~renewal-year fixed effects,
and standard errors are clustered by county (78{,}174 obs.). The key
exclusion: two policyholders mapped into the SFHA close in time but
straddling a calendar-year boundary face different scheduled premium
changes at the same renewal date despite identical map vintages
(Section~5.1.1, p.~18). Identifying assumption: those two groups would
have had the same renewal propensity absent the differential scheduled
price increase.

\paragraph{Headline estimate (Table~1, p.~46).}
\begin{itemize}
\item \textbf{Coefficient on \% premium increase, col.~4 (preferred):}
$\hat\sigma = -0.177$ (s.e.\ 0.040), with state FEs, past county-level
claims controls, and month-written FEs. Coefficients are stable across
columns~1--4 in $[-0.148, -0.177]$.
\item Author phrases this as ``a 1\% increase in premiums decreases policy
renewal by $0.17$ percentage points'' (Section~1, p.~3) and refers to
$|\hat\sigma|=0.177$ as the ``flood insurance price elasticity'' (Section~6.1,
p.~30).
\item \textbf{Mandate adjustment (Appendix~F, p.~74):} using survey
evidence that $\sim$54\% of floodplain buyers carry insurance only because
of the lender mandate, the author scales up to a voluntary-purchaser
elasticity of $0.177/0.46 \approx 0.37$. This is the value plugged into
the welfare model and used to calibrate the isoelastic demand curve $W(I)$
(Section~6.1, p.~30).
\end{itemize}

\paragraph{What kind of object is $\hat\sigma$? (conversion / flag$^\dagger$)}
$\hat\sigma$ is strictly a \emph{semi-elasticity}: $\Delta$(renewal probability,
0--1) per 1\% premium change. The author treats it as a price elasticity
of demand, which assumes the mean renewal rate among the relevant
population is approximately~1. Because newly-mapped policy renewal rates
are very high (the policy data show year-over-year retention well above
85\%; Appendix~E, Figure~A8), the gap between the semi-elasticity
($-0.177$) and a true elasticity ($-0.177/\bar{R}$, with $\bar{R}\approx 0.85$)
is on the order of $0.03$. We therefore record $-0.18$ as the headline
elasticity but flag that:
\begin{itemize}
\item the conversion implicit in the welfare calibration assumes
$\bar{R}\approx 1$ and is mildly attenuating;
\item the mandate-adjusted version ($0.37$) further assumes that
\emph{none} of the observed non-renewals come from mandated buyers and
\emph{all} come from the 46\% voluntary share, which is the strongest
direction the data can support;
\item the local treatment is the response of \emph{policyholders} (those
who already bought) to a known scheduled price path; it is not the
elasticity of new purchase, nor of the broader SFHA stock.
\end{itemize}

\paragraph{Margins captured.}
\textbf{Extensive margin}, renewal only. The panel does not separately
identify cancellations from non-renewals due to home sales. Coverage
amount (intensive margin) is not analyzed in the price-elasticity
exercise.

\paragraph{Threats to identification the author discusses.}
(i) Selection on map vintage within year: addressed by adding state FEs,
past-claims controls, and month-written FEs across columns 2--4 of
Table~1; the coefficient is stable.
(ii) Mandatory purchase requirement: directly addressed via the survey
share adjustment in Appendix~F.
(iii) Salience / information confounds: by design, $\hat\sigma$ is
identified from premium variation \emph{within} the newly-mapped
population, all of whom received the same map-update information signal,
which is what separates this from the floodplain-vs.-non-floodplain
take-up comparison (Table~2). Other dynamics (decaying salience,
mandatory-purchase enforcement) are flagged on p.~18 but cannot be
fully ruled out.

\paragraph{Comparison the author draws.}
The author cites \citet{bradt2021} ($\sim 0.25$) and \citet{wagner2022}
($\sim 0.29$) as similar inelastic estimates and offers the
mandate-adjusted $0.37$ as consistent with these once the mandatory
purchase share is netted out (p.~19, p.~30).

\paragraph{Population to which the estimate applies.}
NFIP policyholders newly mapped into the SFHA between 2015 and 2019, who
purchased at non-floodplain rates and face the scheduled 15\%/yr phase-out
toward floodplain rates. The variation is in scheduled prices set by the
NFIP rate manual, not in market-driven price changes; renewal decisions
are made one year at a time. This is the right population for thinking
about how an existing NFIP policyholder responds when their premium
increment is known and incremental, and a relatively narrow benchmark for
extrapolating to either (a) very large one-shot price increases (cf.\
\citet{collier2026}) or (b) the new-purchase margin.

\paragraph{Other relevant results in the paper (not focal here).}
The author's main object of interest is not $\hat\sigma$ in isolation but
its combination with a positive ``information effect'' that pushes
floodplain take-up \emph{up} despite higher prices. Table~2 (p.~47) shows
floodplain designation raises voluntary take-up by 17~pp conditional on
modeled risk, controls, and demographics ($+21$~pp under an IV using a
zip-code mismapping index). This is a net effect bundling price and
information; the welfare model uses $\hat\sigma$ to net out the price
component and recover an information-only elasticity of demand
(Section~5.2.3, p.~23: holding premiums fixed, floodplain classification
would raise take-up by over 30~pp).

\section{Additional papers (brief synopses, $^\ddagger$)}

\textit{The entries below are based on third-party summaries of the literature
rather than my own read of the source papers. Magnitudes, sample windows, and
specification details should be verified against the originals before citation.
They are included here to broaden the comparison set and triangulate the more
detailed entries above.}

\subsection*{\citet{bradt2021} -- \textit{JEEM}}
A quasi-experimental analysis of the BW-2012 / HFIAA-2014 reforms using NFIP
policy-level data, with the distinctive feature of covering both SFHA properties
(subject to the mandatory purchase requirement) and non-SFHA properties
(voluntary demand). Reports a price elasticity of flood insurance demand of
roughly $-0.25$ to $-0.29$. The SFHA / non-SFHA split is the value added relative
to other BW/HFIAA studies: it lets the reader gauge how much of the elasticity
estimate is dampened by the mandate. Cited as a benchmark by Mulder, Collier
et al., and Hennighausen et al. The paper also includes one of the more careful
tests for adverse selection in flood insurance (none detected), which is why
it is also a standard reference in the welfare literature.

\subsection*{\citet{wagner2022} -- \textit{AEJ: Economic Policy}}
Uses house-level variation in NFIP premiums induced by Biggert--Waters / HFIAA
across 20 Atlantic and Gulf Coast states, 2009--2017, and complements this with
construction-code-induced premium variation as a separate identification source.
Estimates a structural demand model with endogenous adaptation, recovering a
demand elasticity in the $-0.25$ to $-0.30$ range and willingness-to-pay below
$2/3$ of expected loss inside the SFHA. The paper is the natural structural
counterpart to the reduced-form papers in the table and is the most-cited
elasticity estimate in this literature. Wagner is also the source of the
``no adverse selection but adverse selection on adaptation'' finding that
Mulder builds on. House-level (rather than tract- or ZIP-level) premium
variation is what gives Wagner sharper individual-level identification than
the aggregate DiDs in Hennighausen.

\subsection*{\citet{gourevitch2025} -- \textit{J.~Catastrophe Risk and Resilience}}
The first quasi-experimental evidence on the NFIP's Risk Rating 2.0 reform
(2021--2023). Compares ZIP codes by quartile of RR2.0-induced premium increase
and reports an $11$--$39\%$ decline in new NFIP policies and a $5$--$13\%$
decline in renewals, with effects largest in lower-income ZIPs. The summaries
report implied elasticities ranging from $-1.38$ to $-0.26$ across
subgroups$^\dagger$. The headline contribution relative to the BW/HFIAA papers
is twofold: (i) RR2.0 is the current pricing regime, so this is the
forward-looking elasticity, and (ii) the income-quartile heterogeneity is
sharper than anything in the earlier reforms. The wide elasticity range across
subgroups makes a single ``RR2.0 elasticity'' ill-defined; reporting the lower-
and upper-quartile estimates separately is preferable.

\subsection*{\citet{ortega2025} -- \textit{JEEM}}
A companion to \citet{gourevitch2025} on Risk Rating 2.0, but uses individual
NFIP insurance histories (rather than ZIP-level aggregates) and an instrumental
variables strategy. Reports increased exit and reduced entry on the extensive
margin, both inside the SFHA and in its periphery, alongside intensive-margin
adjustments in coverage limits and deductibles --- the latter being something
the ZIP-level papers cannot identify. Available summaries do not give a clean
single elasticity figure; the paper's value-added is the within-policyholder
margin and the SFHA-periphery effect. Worth pairing with Gourevitch et al.\
when discussing the RR2.0 evidence base.

\subsection*{\citet{browne2000} -- \textit{J.~Risk and Uncertainty}}
The classic empirical paper on flood insurance demand. Uses state-year NFIP
aggregates with cross-sectional premium variation and reports two distinct
elasticities: $-0.32$ for policies in force (extensive) and $-1.22$ for total
insured amount (intensive). The intensive-margin elasticity is markedly larger
in magnitude than the extensive-margin one --- a decomposition the recent
quasi-experimental papers (Hennighausen, Collier, Mulder) do not replicate
because their data and designs target the extensive margin. The identification
is weak by modern standards: premiums and policy counts are both equilibrium
objects driven by underlying flood risk and prior loss history, so
state-year correlations are difficult to interpret causally. Useful as a
historical anchor and as the headline reference for the extensive-vs.-intensive
margin distinction.

\subsection*{\citet{atreya2015} -- \textit{Ecological Economics}}
A 32-year county panel for Georgia, with within-county premium variation
identifying take-up responses. Reports elasticities in the $-0.14$ to $-0.31$
range depending on specification. Pre-BW data, so it functions as a
``before-reforms'' benchmark for the later quasi-experimental literature; the
range overlaps with Bradt et al. and Wagner despite very different
identification, which the recent papers cite as a robustness signal.
Identification is weaker than the BW/HFIAA designs (no sharp shock; relies on
within-county time-series variation), so the estimate is best read as a
descriptive long-run aggregate response rather than a causal price elasticity.

\bibliographystyle{abbrvnat}
\begin{thebibliography}{9}
\bibitem[Atreya et al.(2015)]{atreya2015}
Atreya, A., Ferreira, S., Michel-Kerjan, E. (2015).
What drives households to buy flood insurance? New evidence from Georgia.
\textit{Ecological Economics}, 117, 153--161.

\bibitem[Bradt et al.(2021)]{bradt2021}
Bradt, J.\,T., Kousky, C., Wing, O.\,E.\,J. (2021).
Voluntary purchases and adverse selection in the market for flood insurance.
\textit{Journal of Environmental Economics and Management}, 110, 102515.

\bibitem[Browne and Hoyt(2000)]{browne2000}
Browne, M.\,J., Hoyt, R.\,E. (2000).
The demand for flood insurance: Empirical evidence.
\textit{Journal of Risk and Uncertainty}, 20(3), 291--306.

\bibitem[Collier et al.(2026)]{collier2026}
Collier, B.\,L., Huber, T., Jaspersen, J.\,G., Richter, A. (2026).
Homeowners' willingness to insure flood risks as prices increase.
\textit{Working Paper}, April 22, 2026.

\bibitem[Gourevitch et al.(2025)]{gourevitch2025}
Gourevitch, J.\,D., Snyder, S.\,A., Kousky, C. (2025).
Effects of risk-based pricing reform on flood insurance uptake.
\textit{Journal of Catastrophe Risk and Resilience}.

\bibitem[Hennighausen et al.(2023)]{hennighausen2023}
Hennighausen, H., Liao, Y., Nolte, C., Pollack, A. (2023).
Flood insurance reforms, housing market dynamics, and adaptation to climate risks.
\textit{Journal of Housing Economics}, 62, 101953.

\bibitem[Mulder(2024)]{mulder2024}
Mulder, P. (2024).
Mismeasuring risk: The welfare effects of flood risk information.
\textit{Working Paper}, September 24, 2024.

\bibitem[Ortega and Petkov(2025)]{ortega2025}
Ortega, F., Petkov, I. (2025).
To improve is to change: The effects of Risk Rating 2.0 on flood insurance demand.
\textit{Journal of Environmental Economics and Management}.

\bibitem[Wagner(2022)]{wagner2022}
Wagner, K.\,R.\,H. (2022).
Adaptation and adverse selection in markets for natural disaster insurance.
\textit{American Economic Journal: Economic Policy}, 14(3), 380--421.
\end{thebibliography}

\end{document}